% ============================================================
% qp-headfoot.tex —— 页眉页脚模块（全品式导学件，供 main.tex \input）
% 职责：页脚 twoside 奇偶（奇＝右端黑粗小字「章名＋件名」＋页码灰块；偶＝左端页码灰块＋黑粗小字册名）。
%       页眉撤（v4.2，全品 p04-p06 实证无页眉）维持；\qpjieming 宏保留作他件复用参数。
% 参数宏（他件复用改这四条即可）：
%   \qpzhangming 章名　\qpjieming 页眉节名（不排印，保留备他件）　\qpjianming 件名　\qpceming 册名
% v4.3 页脚组（总账C）：
%   页码块 26.2×7.8mm（原 24×7）；数字 ≈11pt（原 10.5）；数字靠版心内侧（原居中）——
%   奇页块贴版心右、数字落块内左（版心）侧；偶页块贴版心左、数字落块内右（版心）侧，
%   内侧留 1.2mm 垫（全品「靠内侧」无精确垫值，出样过目登记）；块几何（底缘 11mm、基线 −2.275
%   ＋raisebox −2.11mm）维持拍板22 口径不变。
%   小字 5.5pt 混重（原 6.5pt 全粗黑）：章名/册名宋体常规＋件名黑体——「混重」按混排字族实现
%   （全品页脚小字实测灰度 58 近黑，色维持黑；具体混重分配出样过目登记）。
%   内容维持我方件名串、不印品牌段（v4.1 已定）。
% 调用关系：main.tex → qp-layout（geometry/颜色）→ 本模块 → qp-titles → qp-blocks。
% ============================================================
\usepackage{fancyhdr}

% ---------- 参数宏（他件复用入口） ----------
\newcommand{\qpzhangming}{第一章\quad 空间向量与立体几何}
\newcommand{\qpjieming}{1.1.1 空间向量及其运算}   % 撤页眉后不排印，保留作他件复用参数
\newcommand{\qpjianming}{导学件}
% v4.4 标定（0908）：册名括号改半角——全角（）是页脚偶页括号墨残留源（断言⑧「全角（）残留」首轮 6 处全出此行，成因登记）
\newcommand{\qpceming}{高中数学\quad 选择性必修第一册(人教B版)}

% ---------- 页码灰块（底 DDDDDD＋黑粗数字，靠版心内侧；v4.3 块 26.2×7.8、数字 11pt） ----------
% 奇页（块贴版心右）：数字落块内左侧（版心侧）；偶页（块贴版心左）：数字落块内右侧。
% H3 片 0909c（E5）：①数字纯黑——\qphuitext 的 \color 是开关、泄漏进页码（旧数字灰 75），数字自带
%   \color{black} 隔离；②数字字面改 \numpnum（NotoSansSC-Medium w500）——全品 p04-p07 复测竖笔
%   5-6px、墨高 39px（2.77mm）；旧 Times Bold 37px/竖笔 3-4px 偏轻，w850 冒 9px 过重（探针七档）；
%   字号 11→10.7pt 使墨高落 39px 档（v1 11pt 实测 40px）；③块出血到纸边——全品复测：块 26.31×7.76、
%   外缘贴纸边（2976/2977px）、内缘距版心 9.0mm（=26.2−17.2）；旧块右缘止于版心。实现＝9mm 宽 \hbox
%   容器：奇页块超出容器右 17.2mm（\hss 吸收）、偶页左伸 17.2mm（前导 \kern），容器净宽 9mm 使文组
%   自然宽度不变（无 overfull）；④数字距块内缘 奇 2.8mm／偶 2.75mm（全品复测 3.03/2.82，去字形承差）
%   ——旧垫 1.2mm 偏小；⑤数字纵向：全品 number ink 底距块底 3.13mm（块底距纸底 11.17）——旧 strut
%   深 2.275 使数字偏低 0.87mm，改 3.145mm（rule 高 7.8 不变）＋页脚 raisebox −2.11→−1.24mm 回位
%   （块底 11.08 维持、数字 ink 底落 14.22≈全品 14.30）。
\newcommand{\qpshuziR}{\hbox to 9mm{\kern-17.2mm\setlength{\fboxsep}{0pt}\colorbox{pnumbg}{%
  \makebox[26.2mm][r]{\rule[-3.145mm]{0pt}{7.8mm}\fontsize{10.7pt}{13pt}\selectfont\color{black}\numpnum\thepage\hspace{2.75mm}}}}}
\newcommand{\qpshuziL}{\hbox to 9mm{\setlength{\fboxsep}{0pt}\colorbox{pnumbg}{%
  \makebox[26.2mm][l]{\rule[-3.145mm]{0pt}{7.8mm}\hspace{2.8mm}\fontsize{10.7pt}{13pt}\selectfont\color{black}\numpnum\thepage}}\hss}}
% v4.3：页脚小字＝5.5pt 混重（章名/册名宋体常规；件名由调用侧挂 \heiti）
% 字替对照 E（0909）：5.5pt→5.24pt（全品字号铁证）＋色黑→灰76（p04 页脚小字暗核众数 76 实测；
%   v4.4 纯黑渲染灰 22 过深，故换 \color{footgray}），混重结构（宋常规＋件名黑体）维持
\definecolor{footgray}{HTML}{4C4C4C}
\newcommand{\qphuitext}{\fontsize{5.24pt}{7pt}\selectfont\color{footgray}}

\pagestyle{fancy}
\fancyhf{}
\renewcommand{\headrulewidth}{0pt}
\fancyfoot[RO]{\raisebox{-1.24mm}[0pt][0pt]{\qphuitext\qpzhangming\quad{\heiti\qpjianming}\hspace{3mm}\qpshuziL}}
\fancyfoot[LE]{\raisebox{-1.24mm}[0pt][0pt]{\qpshuziR\hspace{3mm}\qphuitext\qpceming}}
